% !TEX root = ../main.tex
%==============================================================
%  TÓM TẮT NỘI DUNG ĐỒ ÁN
%==============================================================
\thispagestyle{empty}
\begin{center}
{\Large\bfseries TÓM TẮT NỘI DUNG ĐỒ ÁN}
\end{center}
\vspace{1cm}

Hình ảnh được sinh từ máy tính đóng vai trò quan trọng trong nhiều ứng dụng của
thế giới thực, từ lĩnh vực giải trí như điện ảnh, trò chơi điện tử đến các lĩnh
vực kỹ thuật như kiến trúc và y học. Kết xuất dựa trên vật lý (physically based
rendering) là lớp giải thuật mã hóa các định luật truyền ánh sáng thành mô hình
toán học để tạo ra hình ảnh có độ chân thực cao. Trong số đó, Path Tracing là
phương pháp tổng quát và chính xác, song có chi phí tính toán rất lớn do phải mô
phỏng đường đi của hàng triệu tia sáng bằng tích phân Monte Carlo, dẫn tới thời
gian kết xuất kéo dài và khó đạt tốc độ tương tác.

Trong đồ án này, tôi nghiên cứu và xây dựng một bộ kết xuất Path Tracing chạy
song song trên kiến trúc CUDA của NVIDIA nhằm khai thác sức mạnh của GPU để đạt
tốc độ kết xuất thời gian thực. Hệ thống hiện thực đầy đủ chu trình dò tia trên
GPU, bao gồm: thuật toán Path Tracing với tích lũy theo thời gian và khử răng
cưa; hệ thống vật liệu khuếch tán, kim loại và điện môi; mô hình camera thấu kính
mỏng với hiệu ứng độ sâu trường ảnh và nhòe chuyển động; cơ chế cầu nối
CUDA--OpenGL cho phép hiển thị trực tiếp từ bộ nhớ video; cùng bộ nạp cảnh theo
định dạng Mitsuba~3 và giao diện điều khiển tương tác.

Kết quả cho thấy hệ thống có khả năng kết xuất các cảnh thử nghiệm ở tốc độ
tương tác, đồng thời hội tụ dần về hình ảnh ít nhiễu khi camera giữ nguyên. Đồ án
cũng chỉ ra các hướng tối ưu và mở rộng trong tương lai như cấu trúc tăng tốc BVH
và các kỹ thuật giảm nhiễu nâng cao.

\vspace{1.5cm}
\begin{flushright}
Sinh viên thực hiện\\[0.3cm]
\textit{(Ký và ghi rõ họ tên)}
\end{flushright}
